% Séquence 16 : Proportionnalité (Partie 2)
% Durée estimée : 4 séances
\setseqtitle{Proportionnalité (Partie 2)}
\chapter{Proportionnalité (Partie 2) : Pourcentages et Échelles}
\label{chap:seq16-proportionnalite-2}

\begin{remarkbox}[Rappel de la Partie 1]
	Dans la Partie 1 sur la proportionnalité, nous avons appris :
	\begin{itemize}
		\item Ce qu'est une \textbf{situation de proportionnalité} : deux grandeurs sont proportionnelles si on multiplie toujours par le même nombre pour passer de l'une à l'autre.
		\item Le \textbf{coefficient de proportionnalité} : le nombre constant par lequel on multiplie.
		\item Les \textbf{propriétés de linéarité} : on peut multiplier une colonne ou additionner deux colonnes d'un tableau de proportionnalité.
	\end{itemize}

	\textbf{Dans cette Partie 2}, nous allons voir deux applications importantes de la proportionnalité :
	\begin{itemize}
		\item Les \textbf{pourcentages} (soldes, réductions, statistiques)
		\item Les \textbf{échelles} (cartes, plans, maquettes)
	\end{itemize}
\end{remarkbox}

% ============================================================
% SÉANCE 1 : Les pourcentages
% ============================================================
\seance

\begin{objectifsbox}
	\textbf{Notions :} pourcentages, échelles\\
	\textbf{Compétences :} appliquer un pourcentage, utiliser une échelle (calculer une distance réelle ou sur le plan)\\
	\textbf{Applications :} soldes, statistiques, cartes, plans
\end{objectifsbox}

\section{Comprendre et utiliser les pourcentages}

\subsection{Qu'est-ce qu'un pourcentage ?}

\begin{definitionbox}[Pourcentage]
	Un \textbf{pourcentage} est une fraction dont le dénominateur est égal à 100.
	$$t\% = \dfrac{t}{100}$$
\end{definitionbox}

\begin{examplebox}
	\begin{itemize}
		\item $50\% = \dfrac{50}{100} = \solution{0,5}$ ou $\solution{\dfrac{1}{2}}$ (la moitié)
		\item $25\% = \dfrac{25}{100} = \solution{0,25}$ ou $\solution{\dfrac{1}{4}}$ (le quart)
		\item $10\% = \dfrac{10}{100} = \solution{0,1}$ ou $\solution{\dfrac{1}{10}}$ (le dixième)
		\item $100\% = \dfrac{100}{100} = \solution{1}$ (la totalité, le tout)
	\end{itemize}
\end{examplebox}

\subsection{Appliquer un pourcentage}

\begin{proprietebox}[Règle de calcul]
	Calculer $t\%$ d'un nombre, c'est \textbf{multiplier} ce nombre par $\dfrac{t}{100}$.
\end{proprietebox}

\begin{examplebox}[Exemple guidé]
	Dans une classe de 25 élèves, $60\%$ sont des demi-pensionnaires. Combien cela représente-t-il d'élèves ?

	\textbf{Calcul :}
	$$25 \times 60\% = 25 \times \dfrac{60}{100}$$

	On peut calculer de plusieurs manières :
	\begin{itemize}
		\item \textbf{Méthode 1 :} Convertir le pourcentage en nombre décimal : $25 \times 0,60 = \solution{15}$
		\item \textbf{Méthode 2 :} Multiplier d'abord par 60, puis diviser par 100 : $(25 \times 60) \div 100 = 1~500 \div 100 = \solution{15}$
		\item \textbf{Méthode 3 :} Diviser d'abord par 100, puis multiplier par 60 : $(25 \div 100) \times 60 = 0,25 \times 60 = \solution{15}$
	\end{itemize}

	\textbf{Réponse :} Il y a \solution{15} élèves demi-pensionnaires dans cette classe.
\end{examplebox}

\begin{activitybox}[Les soldes]
	Un pantalon coûte 40~€. Il est soldé à $-20\%$.
	\begin{enumerate}
		\item Quel est le montant de la réduction ?

		\solution{
			\textbf{Calcul :} Pour trouver le montant de la réduction, on calcule 20\% de 40~€ :
			$$40 \times \dfrac{20}{100} = 40 \times 0,2 = \solution{8}~\text{€}$$

			\textbf{Réponse :} La réduction est de 8~€.
		}

		\item Quel est le nouveau prix ?

		\solution{
			\textbf{Calcul :} On soustrait la réduction du prix initial :
			$$40 - 8 = \solution{32}~\text{€}$$

			\textbf{Réponse :} Le nouveau prix après réduction est de 32~€.
		}
	\end{enumerate}
\end{activitybox}

\subsection{Calculer un pourcentage à partir de deux quantités}

Parfois, on connaît la partie et le tout, et on veut trouver le pourcentage.

\begin{methodebox}[Méthode pour trouver le pourcentage]
	Pour trouver quel pourcentage représente une quantité par rapport à une autre :
	\begin{enumerate}
		\item On divise la partie par le tout
		\item On multiplie le résultat par 100
	\end{enumerate}
	$$\text{Pourcentage} = \dfrac{\text{Partie}}{\text{Tout}} \times 100$$
\end{methodebox}

\begin{examplebox}[Application]
	Dans une classe de 25 élèves, 15 sont des filles. Quel pourcentage de la classe représentent les filles ?

	\textbf{Calcul :} On applique la formule $\dfrac{\text{Partie}}{\text{Tout}} \times 100$ :
	$$\dfrac{15}{25} \times 100 = 0,6 \times 100 = \solution{60}\%$$

	\textbf{Détail :}
	\begin{itemize}
		\item On divise d'abord : $15 \div 25 = 0,6$
		\item Puis on multiplie par 100 : $0,6 \times 100 = 60$
	\end{itemize}

	\textbf{Réponse :} Les filles représentent \solution{60}\% de la classe.
\end{examplebox}

\subsection{Pourcentages particuliers et calcul mental}

Certains pourcentages correspondent à des fractions simples qu'il est utile de connaître pour calculer rapidement.

\begin{summarybox}[À retenir par cœur]
	\begin{itemize}
		\item Prendre \textbf{50\%}, c'est prendre la \textbf{moitié} (diviser par 2).
		\item Prendre \textbf{25\%}, c'est prendre le \textbf{quart} (diviser par 4).
		\item Prendre \textbf{75\%}, c'est prendre les \textbf{trois-quarts} ($\times 3$ puis $\div 4$).
		\item Prendre \textbf{10\%}, c'est prendre le \textbf{dixième} (diviser par 10).
	\end{itemize}
\end{summarybox}

% ============================================================
% SÉANCE 2 : Augmentations et réductions
% ============================================================
\seance

\section{Augmentations et réductions en pourcentage}

Dans la vie quotidienne, on rencontre souvent des situations où une quantité augmente ou diminue d'un certain pourcentage.

\subsection{Appliquer une réduction}

\begin{methodebox}[Méthode pour une réduction]
	Pour calculer le nouveau prix après une réduction de $t\%$ :
	\begin{enumerate}
		\item Calculer le montant de la réduction : Prix initial $\times \dfrac{t}{100}$
		\item Soustraire la réduction du prix initial
	\end{enumerate}

	\textbf{Ou en une seule étape :}
	$$\text{Nouveau prix} = \text{Prix initial} \times \left(1 - \dfrac{t}{100}\right)$$
\end{methodebox}

\begin{examplebox}[Soldes : méthode complète]
	Un ordinateur coûte 600~€. Il a une réduction de $30\%$.

	\textbf{Méthode 1 (en deux étapes) :}
	\begin{itemize}
		\item Réduction : $600 \times 0,30 = \solution{180}$~€
		\item Nouveau prix : $600 - 180 = \solution{420}$~€
	\end{itemize}

	\textbf{Méthode 2 (en une étape) :}
	\begin{itemize}
		\item Une réduction de $30\%$ signifie qu'on garde $100\% - 30\% = 70\%$ du prix.
		\item Nouveau prix : $600 \times 0,70 = \solution{420}$~€
	\end{itemize}
\end{examplebox}

\subsection{Appliquer une augmentation}

\begin{methodebox}[Méthode pour une augmentation]
	Pour calculer la nouvelle valeur après une augmentation de $t\%$ :
	$$\text{Nouvelle valeur} = \text{Valeur initiale} \times \left(1 + \dfrac{t}{100}\right)$$
\end{methodebox}

\begin{examplebox}[Augmentation de prix]
	Le prix du pain est de 1,20~€. Il augmente de $5\%$. Quel est le nouveau prix ?

	\textbf{Calcul :}
	\begin{itemize}
		\item Une augmentation de $5\%$ signifie qu'on a $100\% + 5\% = 105\%$ du prix initial.
		\item Nouveau prix : $1,20 \times 1,05 = \solution{1,26}$~€
	\end{itemize}
\end{examplebox}

\subsection{Situations avec pourcentages successifs}

\begin{remarkbox}[Attention aux pourcentages successifs !]
	Attention : $-20\%$ puis $-10\%$ n'est \textbf{pas égal} à $-30\%$ !

	On applique les pourcentages l'un après l'autre, sur le résultat précédent.
\end{remarkbox}

\begin{examplebox}[Soldes successifs]
	Un article coûte 100~€. Il est soldé à $-20\%$, puis il y a encore $-10\%$ sur le prix soldé.

	\textbf{Calcul :}
	\begin{itemize}
		\item Après la première réduction : $100 \times 0,80 = \solution{80}$~€
		\item Après la deuxième réduction : $80 \times 0,90 = \solution{72}$~€
	\end{itemize}

	\textbf{Vérification :} Si c'était $-30\%$, on aurait : $100 \times 0,70 = 70$~€ $\neq 72$~€
\end{examplebox}

% ============================================================
% SÉANCE 3 : Les échelles
% ============================================================
\seance

\section{Les échelles}

L'échelle d'un plan ou d'une carte est un cas particulier de proportionnalité.

\begin{definitionbox}[Échelle]
	L'\textbf{échelle} d'une carte (ou d'un plan) est le coefficient de proportionnalité permettant de passer des distances réelles aux distances sur la carte, exprimées dans la \textbf{même unité}.

	$$\text{Échelle} = \dfrac{\text{Distance sur le plan}}{\text{Distance réelle}}$$
\end{definitionbox}

\begin{remarkbox}
	Une échelle est souvent notée sous forme de fraction $\dfrac{1}{n}$.
	Cela signifie que \textbf{1 cm sur le plan} représente \textbf{n cm en réalité}.
\end{remarkbox}

\begin{examplebox}[Comprendre une échelle]
	Une carte est à l'échelle $\dfrac{1}{200~000}$.
	\begin{itemize}
		\item Cela signifie que \textbf{1 cm} sur la carte représente \textbf{200 000 cm} en réalité.
		\item On convertit souvent la distance réelle en mètres ou kilomètres pour que ce soit plus parlant :
		\begin{itemize}
			\item $200~000$ cm $= 200~000 \div 100 = 2~000$ m (car 1 m = 100 cm)
			\item $2~000$ m $= 2~000 \div 1~000 = \solution{2}$ km (car 1 km = 1 000 m)
		\end{itemize}
		\item Donc 1 cm sur la carte représente \solution{2} km en réalité.
	\end{itemize}
\end{examplebox}

\begin{methodebox}[Utiliser un tableau de proportionnalité]
	Pour résoudre des problèmes d'échelle, on utilise un tableau de proportionnalité :
	\begin{center}
		\renewcommand{\arraystretch}{1.5}
		\begin{tabular}{|c|c|c|}
			\hline
			\textbf{Distance sur le plan (cm)} & 1 & $d$ \\
			\hline
			\textbf{Distance réelle (cm)} & $n$ & $D$ \\
			\hline
		\end{tabular}
	\end{center}
	\textbf{Attention à bien mettre les distances réelles en cm dans le tableau !}
\end{methodebox}

\begin{exercisebox}[Application 1 : du plan à la réalité]
	Sur un plan à l'échelle $\dfrac{1}{500}$, la chambre mesure 4 cm de long.

	\textbf{Mots pour comprendre :} $1$ cm sur le plan $\rightarrow$ $500$ cm en vrai.

	\solution{
		\textbf{Calcul :}
		\begin{itemize}
			\item Longueur réelle en cm : La chambre mesure 4 cm sur le plan. On multiplie par 500 :
			$$4 \times 500 = \solution{2~000}~\text{cm}$$

			\item Longueur réelle en m : On convertit 2 000 cm en mètres (1 m = 100 cm) :
			$$2~000 \div 100 = \solution{20}~\text{m}$$
		\end{itemize}

		\textbf{Réponse :} La chambre mesure 20 m de long en réalité.
	}
\end{exercisebox}

\subsection{Calculer une distance sur le plan (sens inverse)}

Parfois, on connaît la distance réelle et on veut trouver la distance sur le plan.

\begin{methodebox}[De la réalité au plan]
	Pour calculer une distance sur le plan à partir d'une distance réelle :
	\begin{enumerate}
		\item Convertir la distance réelle en cm
		\item Diviser par le coefficient de l'échelle
	\end{enumerate}

	Si l'échelle est $\dfrac{1}{n}$ :
	$$\text{Distance sur le plan} = \dfrac{\text{Distance réelle (en cm)}}{n}$$
\end{methodebox}

\begin{examplebox}[De la réalité au plan]
	On veut représenter une route de 3 km sur une carte à l'échelle $\dfrac{1}{100~000}$. Quelle sera la longueur sur la carte ?

	\textbf{Étape 1 :} Convertir 3 km en cm
	\begin{itemize}
		\item $3$ km $= 3 \times 1~000 = 3~000$ m (car 1 km = 1 000 m)
		\item $3~000$ m $= 3~000 \times 100 = \solution{300~000}$ cm (car 1 m = 100 cm)
	\end{itemize}

	\textbf{Étape 2 :} Diviser par le coefficient d'échelle
	$$\text{Distance sur la carte} = \dfrac{300~000}{100~000} = \solution{3}~\text{cm}$$

	\textbf{Vérification :} L'échelle $\dfrac{1}{100~000}$ signifie que 1 cm sur la carte = 100 000 cm en réalité. Donc 3 cm sur la carte représentent bien $3 \times 100~000 = 300~000$ cm = 3 km.

	\textbf{Réponse :} La route mesurera \solution{3} cm sur la carte.
\end{examplebox}

\begin{exercisebox}[Application]
	Sur une carte à l'échelle $\dfrac{1}{25~000}$, deux villes sont distantes de 8 cm.

	\begin{enumerate}
		\item Quelle est la distance réelle en cm ?

		\solution[5cm]{
			\textbf{Calcul :} L'échelle $\dfrac{1}{25~000}$ signifie que 1 cm sur la carte = 25 000 cm en réalité.

			Pour 8 cm sur la carte : $8 \times 25~000 = \solution{200~000}$ cm

			\textbf{Réponse :} La distance réelle est de 200 000 cm.
		}

		\item Quelle est la distance réelle en km ?

		\solution[5cm]{
			\textbf{Calcul :} On convertit 200 000 cm en kilomètres :
			\begin{itemize}
				\item $200~000$ cm $= 200~000 \div 100 = 2~000$ m
				\item $2~000$ m $= 2~000 \div 1~000 = \solution{2}$ km
			\end{itemize}

			\textbf{Réponse :} La distance réelle entre les deux villes est de 2 km.
		}
	\end{enumerate}
\end{exercisebox}

\begin{remarkbox}[Pour s'entra\^{i}ner]
	\textbf{Cahier Transmath 6e - Chapitre 16 :}
	\begin{itemize}
		\item Fiche 118 (p.129) \og \'Echelles\fg{} : ex. 51 \`a 56
	\end{itemize}
\end{remarkbox}

% ============================================================
% SÉANCE 4 : Bilan de la séquence
% ============================================================
\seance

\begin{remarkbox}[Pour s'entra\^{i}ner]
	\textbf{Cahier Transmath 6e - Chapitre 16 :}
	\begin{itemize}
		\item Fiche 117 (p.128) \og Application d'un pourcentage\fg{} : ex. 42 \`a 50
		\item Fiche 118 (p.129) \og \'Echelles\fg{} : ex. 51 \`a 56
		\item Fiche 119 (p.130) \og R\'esoudre des probl\`emes\fg{} : probl\`emes de synth\`ese
	\end{itemize}
\end{remarkbox}

\section{Bilan de la séquence}

\begin{bilanbox}
	\textbf{Ce que je dois savoir :}
	\begin{itemize}
		\item[$\checkmark$] Savoir qu'un pourcentage est une fraction sur 100
		\item[$\checkmark$] Connaître les pourcentages usuels (50\%, 25\%, 10\%, 75\%)
		\item[$\checkmark$] Savoir ce qu'est une échelle et comment elle fonctionne
		\item[$\checkmark$] Connaître les formules d'augmentation et de réduction
	\end{itemize}

	\textbf{Ce que je dois savoir faire :}
	\begin{itemize}
		\item[$\checkmark$] Calculer un pourcentage d'un nombre
		\item[$\checkmark$] Trouver quel pourcentage représente une quantité par rapport à une autre
		\item[$\checkmark$] Appliquer une augmentation ou une réduction en pourcentage
		\item[$\checkmark$] Calculer une distance réelle à partir d'un plan (et inversement)
		\item[$\checkmark$] Utiliser un tableau de proportionnalité et convertir les unités
	\end{itemize}
\end{bilanbox}

\begin{summarybox}[L'essentiel]
	\textbf{Pourcentages :}
	\begin{itemize}
		\item $t\% = \dfrac{t}{100}$
		\item Calculer $t\%$ de $N$ : $N \times \dfrac{t}{100}$
		\item Trouver un pourcentage : $\dfrac{\text{Partie}}{\text{Tout}} \times 100$
		\item Réduction de $t\%$ : multiplier par $(1 - \dfrac{t}{100})$
		\item Augmentation de $t\%$ : multiplier par $(1 + \dfrac{t}{100})$
		\item Astuces : $50\% = \div 2$ ; $25\% = \div 4$ ; $10\% = \div 10$ ; $75\% = \times 3$ puis $\div 4$
	\end{itemize}

	\textbf{Échelles :}
	\begin{itemize}
		\item Échelle $= \dfrac{\text{Distance sur le plan}}{\text{Distance réelle}}$ (mêmes unités !)
		\item Échelle $\dfrac{1}{n}$ : 1 cm sur le plan = $n$ cm en réalité
		\item Plan $\rightarrow$ Réalité : multiplier par $n$
		\item Réalité $\rightarrow$ Plan : diviser par $n$
		\item \textbf{Toujours} convertir en cm avant de faire les calculs !
	\end{itemize}
\end{summarybox}
